\section{Introduction}

In medical intervention research including in particular randomized controlled trials (RCTs) there is a trend towards increased use of patient reported outcomes; quality of life scores being one of the most prominent of these. Established statistical practice is to compare these between treatment groups using non-parametric tests such as the Mann--Whitney--Wilcoxon rank-sum test \citep{wilcoxon3001968individual,mann1947test}, henceforth Wilcoxon, since distributions are rarely normal or symmetric \textbf{[TODO: References needed to current statistical practice]}. When all patients are alive and able and willing to provide quality of life scores at the scheduled measurement time this procedure works very well. Not all clinical settings, however, will have all participants alive at time of scheduled assessment of quality of life. This is particularly true in trials within intensive care where mortality of approximately 30\% is not uncommon. This is the setting that motivated this work. A pragmatic solution which is gaining popularity is to simply give the deceased patients the lowest possible score and then use a Wilcoxon test or similar to compare the groups \textbf{[TODO: References needed]}. This paper demonstrates that this approach can lead to a dramatic loss of statistical power, and we will introduce an alternative statistical test statistic that averts this power loss. Our proposed test procedure is implemented in an R package and made publicly available.

The key insight in the proposed test is to incorporate that the outcome is actually a two-dimensional outcome of a very special type, and that the constructed combined outcome follows a continuous-singular mixture distribution. This unusual distribution is why one cannot resort to non-parametric Wilcoxon tests since the singular component in the distribution of the combined outcome will get reduced to simple ties. It is noted that the handling of ties in standard statistical software varies and is opaque. However, the handling of ties is not the main reason why the Wilcoxon test suffers power loss. The main reason is that the null hypothesis in these Wilcoxon type tests (stochastic domination) does not handle the empirical fact that treatments might influence mortality and quality of life differently.

As an alternative we propose to model the binary component (i.e., survival) and the continuous part (i.e., actual quality of life) separately but to construct a single test for no treatment effect on either based on a likelihood ratio. We can thus provide a single p-value for the hypothesis of no treatment effect on the extended quality of life where death is given the lowest possible score. To accommodate potential non-normality of the recorded quality of life scores we consider both a parametric and a semi-parametric version of the test where the semi-parametric version is as widely applicable as the Wilcoxon test. Both test procedures will provide effect estimates of mean differences between treatment groups based on the combined outcome along with confidence intervals. This is also an added benefit compared to the Wilcoxon-type testing approach.

It should be noted that while we in this paper exemplify the procedure using quality of life and mortality the method is applicable in any setting where a single value of a combined outcome has a special interpretation compared to an otherwise continuous scale. Another example from intensive care research consider the outcome ``days alive and out of hospital within 90 days from randomization''. Here in-hospital fatalities will all have the value 0, while everybody else will have outcomes ranging from 0 to 90. Such outcomes are also routinely analyzed using Wilcoxon-type tests. Our proposed test would increase power while also providing a mean effect estimate along with confidence interval. Note also that the developed mathematical method allows straightforwardly for the inclusion of confounding variables or other adjustment factors. The method itself is therefore just as applicable in non-randomized studies or epidemiological studies in general.

The rest of the paper is structured as follows. The next section introduces the mathematical setup as well as our novel test procedure. Section~\ref{sec:software-implementation} describes the R implementation, and Section~\ref{sec:simulation-study} presents a simulation study illustrating the substantial power gains. Section~\ref{sec:application} re-analyzes a COVID-19 trial from 2021, and Section~\ref{sec:discussion} concludes with a discussion. Additional simulation results and mathematical proofs are contained in Appendix~\ref{app:supplementary-simulation-results} and Appendix~\ref{app:supplementary-derivation}.
